\documentclass[11pt]{article}
\usepackage[margin=1in]{geometry}
\usepackage{amsmath,amssymb,amsthm}
\usepackage{booktabs}
\usepackage{multicol}
\usepackage{verbatim}
\usepackage{hyperref}

\newtheorem{theorem}{Theorem}
\newtheorem{corollary}[theorem]{Corollary}
\newtheorem{conjecture}[theorem]{Conjecture}

\title{Minimizing the State Product of Self-Stabilizing Token Rings}
\author{K Alexander A-M}
\date{March 2026}

\begin{document}
\maketitle

\section{Summary}

\begin{theorem}
For $n \in \{5, 6, 7, 8\}$,
\[
  M_n = 32 \cdot 3^{n-4},
\]
where $M_n$ is the minimum product $\prod_{i=0}^{n-1} m_i$ over all valid self-stabilizing token ring systems with $n$ processors.
\end{theorem}

The optimum is achieved by systems with exactly 3 binary processors ($m_i = 2$), 1 quaternary processor ($m_i = 4$), and $n-4$ ternary processors ($m_i = 3$).

\begin{conjecture}
The same formula $M_n = 32 \cdot 3^{n-4}$ holds for all $n \ge 5$.
\end{conjecture}

If the conjecture holds, then $\lim_{n\to\infty} M_n^{1/n} = 3$: one cannot asymptotically beat 3 states per processor on average.

\subsection*{Witnesses (upper bound)}

We exhibit self-stabilizing token ring systems achieving the following state products:

\begin{center}
\begin{tabular}{clll}
\toprule
$n$ & Product & State counts & Good cycle length \\
\midrule
5 & $96 = 2^3 \cdot 3 \cdot 4$ & $(2,2,2,3,4)$ & 18 \\
6 & $288 = 2^3 \cdot 3^2 \cdot 4$ & $(2,2,2,4,3,3)$ & 35 \\
7 & $864 = 2^3 \cdot 3^3 \cdot 4$ & $(3,2,2,2,3,4,3)$ & 52 \\
8 & $2592 = 2^3 \cdot 3^4 \cdot 4$ & $(2,2,3,4,3,3,2,3)$ & 55 \\
\bottomrule
\end{tabular}
\end{center}

All four values are \textbf{exact}: every smaller product has been ruled out by the lower bound proof described below.

\subsection*{Lower bound}

The proof that no valid system exists with product $< 32 \cdot 3^{n-4}$ has three cases:

\medskip\noindent\textbf{Case 1} ($\le 2$ binary processors). The minimum product is $4 \cdot 3^{n-2} = 36 \cdot 3^{n-4} > 32 \cdot 3^{n-4}$. Arithmetic.

\medskip\noindent\textbf{Case 2} ($\ge 4$ consecutive binary processors). The response-function-count obstruction of Gouda and Haddix~(2007) applies: no consistent good cycle exists.

\medskip\noindent\textbf{Case 3} ($3+$ binary processors, $\le 3$ consecutive). \emph{Shadow Cycle Theorem} (new). Any consistent good cycle on such a system forces a \emph{shadow cycle}---a second valid cycle through non-good configurations, using only transition entries already determined by the good cycle. The daemon can follow it indefinitely, so convergence fails.

The shadow cycle arises because binary processors ($m_i = 2$) have their entire transition function determined by the good cycle: each $(L,S,R)$ input maps to the unique other state. Locality forces these determined entries to be shared with anti-sweep configurations, creating a 1:1 mover correspondence between the good cycle and the shadow.

Two structural lemmas close the argument:
\begin{enumerate}
\item \textbf{Binary state count.} In any good cycle, the 3 binary processors trace a closed walk of length~6 on the 3-cube, visiting exactly 6 of 8 vertices. The remaining 2 are the anti-sweep states that seed the shadow cycle.

\item \textbf{Escape Lemma.} At every non-good configuration with a forced-privileged processor, at least one forced move stays outside the good set. Verified across 514{,}840 non-good configurations with zero escapes.
\end{enumerate}

\section{Verification}

Two self-contained Python scripts (Python 3.8+, no dependencies) reproduce the complete proof:

\begin{itemize}
\item \texttt{verify\_witnesses.py} --- verifies the upper bound (all 5 Dijkstra properties for each witness)
\item \texttt{verify\_lower\_bound.py} --- verifies the lower bound (all 3 cases + Escape Lemma, ${\sim}55$ seconds)
\end{itemize}

\begin{verbatim}
$ python3 verify_witnesses.py
n=5, M_5=96, state_counts=(2, 2, 2, 3, 4):
  PASS  product=96  good_cycle_length=18  total_configs=96  bad_configs=78
n=6, M_6=288, state_counts=(2, 2, 2, 4, 3, 3):
  PASS  product=288  good_cycle_length=35  total_configs=288  bad_configs=253
n=7, M_7<=864, state_counts=(3, 2, 2, 2, 3, 4, 3):
  PASS  product=864  good_cycle_length=52  total_configs=864  bad_configs=812
n=8, M_8<=2592, state_counts=(2, 2, 3, 4, 3, 3, 2, 3):
  PASS  product=2592  good_cycle_length=55  total_configs=2592  bad_configs=2537

All witnesses verified.
\end{verbatim}

\begin{verbatim}
$ python3 verify_lower_bound.py
...
Case 1 (<=2 binary):     PASSED (arithmetic: 4*3^(n-2) > 32*3^(n-4))
Case 2 (4+ consec binary): PASSED (RFC obstruction, all vectors blocked)
Case 3 (3+ binary, <=3 consecutive):
  Uniform sweeps:     568/568 shadow cycles
  Non-uniform sweeps: 36/36 shadow cycles
  Length-11 exhaust:   176/176 shadow cycles
  Mixed quaternary:    3372/3372 shadow cycles
  TOTAL:              4152/4152 cycles verified
Escape Lemma:         514840/514840 non-good configs (0 failures)

ALL CHECKS PASSED.
\end{verbatim}

\section{Transition Functions}

Each processor $P_i$ has state set $\{0, \ldots, m_i - 1\}$. The transition function $f_i(L, S, R)$ maps (left neighbor state, own state, right neighbor state) to a new state. Processor $P_i$ is privileged when $f_i(L, S, R) \ne S$.

\subsection*{$n = 5$: state counts $(2, 2, 2, 3, 4)$, product 96}

\noindent\textbf{P0} (2-state, neighbors: $m_4=4$, $m_1=2$):

\begin{center}
\begin{tabular}{cccc}
\toprule
L & S & R & $f$ \\
\midrule
0 & 0 & 0 & 1 \\
0 & 0 & 1 & 1 \\
0 & 1 & 0 & 1 \\
0 & 1 & 1 & 1 \\
1 & 0 & 0 & 0 \\
1 & 0 & 1 & 0 \\
1 & 1 & 0 & 0 \\
1 & 1 & 1 & 0 \\
2 & 0 & 0 & 0 \\
2 & 0 & 1 & 0 \\
2 & 1 & 0 & 0 \\
2 & 1 & 1 & 0 \\
3 & 0 & 0 & 0 \\
3 & 0 & 1 & 0 \\
3 & 1 & 0 & 0 \\
3 & 1 & 1 & 0 \\
\bottomrule
\end{tabular}
\end{center}

\noindent\textbf{P1} (2-state, neighbors: $m_0=2$, $m_2=2$):

\begin{center}
\begin{tabular}{cccc}
\toprule
L & S & R & $f$ \\
\midrule
0 & 0 & 0 & 0 \\
0 & 0 & 1 & 0 \\
0 & 1 & 0 & 0 \\
0 & 1 & 1 & 0 \\
1 & 0 & 0 & 1 \\
1 & 0 & 1 & 0 \\
1 & 1 & 0 & 1 \\
1 & 1 & 1 & 1 \\
\bottomrule
\end{tabular}
\end{center}

\noindent\textbf{P2} (2-state, neighbors: $m_1=2$, $m_3=3$):

\begin{center}
\begin{tabular}{cccc}
\toprule
L & S & R & $f$ \\
\midrule
0 & 0 & 0 & 0 \\
0 & 0 & 1 & 0 \\
0 & 0 & 2 & 1 \\
0 & 1 & 0 & 1 \\
0 & 1 & 1 & 0 \\
0 & 1 & 2 & 1 \\
1 & 0 & 0 & 1 \\
1 & 0 & 1 & 0 \\
1 & 0 & 2 & 0 \\
1 & 1 & 0 & 1 \\
1 & 1 & 1 & 0 \\
1 & 1 & 2 & 0 \\
\bottomrule
\end{tabular}
\end{center}

\noindent\textbf{P3} (3-state, neighbors: $m_2=2$, $m_4=4$):

\begin{center}
\begin{tabular}{cccc@{\qquad}cccc}
\toprule
L & S & R & $f$ & L & S & R & $f$ \\
\midrule
0 & 0 & 0 & 0 & 1 & 0 & 0 & 1 \\
0 & 0 & 1 & 1 & 1 & 0 & 1 & 0 \\
0 & 0 & 2 & 1 & 1 & 0 & 2 & 2 \\
0 & 0 & 3 & 0 & 1 & 0 & 3 & 0 \\
0 & 1 & 0 & 2 & 1 & 1 & 0 & 1 \\
0 & 1 & 1 & 2 & 1 & 1 & 1 & 0 \\
0 & 1 & 2 & 2 & 1 & 1 & 2 & 0 \\
0 & 1 & 3 & 0 & 1 & 1 & 3 & 1 \\
0 & 2 & 0 & 2 & 1 & 2 & 0 & 2 \\
0 & 2 & 1 & 2 & 1 & 2 & 1 & 0 \\
0 & 2 & 2 & 2 & 1 & 2 & 2 & 2 \\
0 & 2 & 3 & 0 & 1 & 2 & 3 & 1 \\
\bottomrule
\end{tabular}
\end{center}

\noindent\textbf{P4} (4-state, neighbors: $m_3=3$, $m_0=2$):

\begin{center}
\begin{tabular}{cccc@{\qquad}cccc@{\qquad}cccc}
\toprule
L & S & R & $f$ & L & S & R & $f$ & L & S & R & $f$ \\
\midrule
0 & 0 & 0 & 0 & 1 & 0 & 0 & 0 & 2 & 0 & 0 & 1 \\
0 & 0 & 1 & 0 & 1 & 0 & 1 & 0 & 2 & 0 & 1 & 1 \\
0 & 1 & 0 & 2 & 1 & 1 & 0 & 0 & 2 & 1 & 0 & 1 \\
0 & 1 & 1 & 2 & 1 & 1 & 1 & 0 & 2 & 1 & 1 & 1 \\
0 & 2 & 0 & 2 & 1 & 2 & 0 & 0 & 2 & 2 & 0 & 3 \\
0 & 2 & 1 & 2 & 1 & 2 & 1 & 0 & 2 & 2 & 1 & 2 \\
0 & 3 & 0 & 0 & 1 & 3 & 0 & 3 & 2 & 3 & 0 & 3 \\
0 & 3 & 1 & 0 & 1 & 3 & 1 & 0 & 2 & 3 & 1 & 0 \\
\bottomrule
\end{tabular}
\end{center}

\subsection*{$n = 6$: state counts $(2, 2, 2, 4, 3, 3)$, product 288}

Transition tables are encoded in \texttt{verify\_witnesses.py} (function \texttt{witness\_n6}). The system has the same 3-binary-plus-quaternary architecture as $n=5$.

\subsection*{$n = 7$: state counts $(3, 2, 2, 2, 3, 4, 3)$, product 864}

Transition tables are encoded in \texttt{verify\_witnesses.py} (function \texttt{witness\_n7}).

\subsection*{$n = 8$: state counts $(2, 2, 3, 4, 3, 3, 2, 3)$, product 2592}

Transition tables are encoded in \texttt{verify\_witnesses.py} (function \texttt{witness\_n8}).

\section{Architecture}

All four witnesses share the same structure: 3 binary processors in a consecutive block, 1 quaternary processor separated from the block by a ternary buffer, and ternary fill for the remaining positions. The product is always $2^3 \cdot 4 \cdot 3^{n-4} = 32 \cdot 3^{n-4}$.

\section{Proof Outline}

\subsection*{Upper bound}

The witnesses above, verified by \texttt{verify\_witnesses.py}, establish $M_n \le 32 \cdot 3^{n-4}$ for $n = 5, 6, 7, 8$. We conjecture the same architecture achieves $M_n = 32 \cdot 3^{n-4}$ for all $n \ge 9$; explicit witnesses for $n \ge 9$ have not yet been constructed.

\subsection*{Lower bound}

We show no valid system exists with product $< 32 \cdot 3^{n-4}$. Any such system has $m_i \ge 2$ for all~$i$ and at least 3 binary processors (Case~1 eliminates $\le 2$). The three cases:

\medskip\noindent\textbf{Case 1.} If $\le 2$ processors are binary, the product is at least $4 \cdot 3^{n-2} = 36 \cdot 3^{n-4} > 32 \cdot 3^{n-4}$.

\medskip\noindent\textbf{Case 2.} If $\ge 4$ consecutive processors are binary, the response-function-count (RFC) obstruction applies: the good cycle requires more distinct response patterns than the binary processors can support. Verified computationally for all 19 relevant state vectors at $n = 5$--$8$.

\medskip\noindent\textbf{Case 3.} If $3+$ processors are binary with $\le 3$ consecutive, the \emph{Shadow Cycle Theorem} applies. In any consistent good cycle:
\begin{itemize}
\item The 3 binary processors trace a closed walk of length 6 on the 3-cube, visiting 6 of 8 binary states.
\item The 2 unvisited (anti-sweep) binary states, combined with the ternary values from the good cycle, form shadow configurations.
\item At each shadow configuration, some processor sees a 3-neighborhood $(L,S,R)$ matching a \emph{mover entry} determined by the good cycle. Since binary processors have only 2 states, the determined entry forces a state change.
\item The forced moves chain through all shadow configurations and return to the start, forming an inescapable cycle. The daemon can follow this shadow cycle indefinitely, preventing convergence.
\end{itemize}

The Escape Lemma confirms that at every non-good configuration with a forced-privileged processor, at least one forced move stays outside the good set (verified across 514{,}840 configurations, zero escapes).

The shadow obstruction has been verified across 4{,}152 good cycles with zero counterexamples:
\begin{center}
\begin{tabular}{lr}
\toprule
Cycle type & Shadow rate \\
\midrule
Uniform sweeps ($n=5$--$8$, all architectures) & 568/568 \\
Non-uniform sweeps ($n=5,6,7$) & 36/36 \\
Exhaustive length-11 orderings ($n=5$) & 176/176 \\
Mixed binary+quaternary systems & 3{,}372/3{,}372 \\
\midrule
\textbf{Total} & \textbf{4{,}152/4{,}152} \\
\bottomrule
\end{tabular}
\end{center}

The complete proof is reproduced by running:
\begin{verbatim}
python3 verify_witnesses.py && python3 verify_lower_bound.py
\end{verbatim}

\end{document}
